\documentclass[11pt]{article}
\usepackage{amsmath,amssymb,amsthm}
\usepackage{algorithm}
\usepackage{algorithmic}
\usepackage{enumitem}
\usepackage{hyperref}
\usepackage{geometry}
\geometry{margin=1in}
\newtheorem{lemma}{Lemma}
\newtheorem{theorem}{Theorem}
\newcommand{\E}{\mathbb{E}}
\newcommand{\R}{\mathbb{R}}

\begin{document}

\section*{Local SGD (Federated Averaging)}

\section*{Mathematical Formulation}
Consider $K$ clients indexed by $k\in\{1,\dots,K\}$.  
Each client possesses a local convex smooth loss $f_k:\R^d\rightarrow\R$ and a stochastic gradient oracle
\[
g_t^{(k)}:=\nabla f_k\big(w_t^{(k)};\,\xi_t^{(k)}\big),
\]
where $\xi_t^{(k)}$ is a random data sample drawn i.i.d.\ from the local distribution of client $k$.

Let $E\in\mathbb{N}$ denote the number of local SGD steps performed between two communication rounds and let $\eta>0$ be a (possibly time‑varying) learning rate.  
Denote by $w_t$ the global model after the $t$‑th communication.  
The algorithm proceeds as follows.

\medskip
\noindent\textbf{Local updates (for each client $k$).}  
For $e=0,1,\dots,E-1$,
\begin{equation}
\label{eq:local}
w_{t,e+1}^{(k)} \;=\; w_{t,e}^{(k)} \;-\; \eta_t\, g_{t,e}^{(k)},
\qquad g_{t,e}^{(k)}:=\nabla f_k\big(w_{t,e}^{(k)};\,\xi_{t,e}^{(k)}\big),
\end{equation}
with initialization $w_{t,0}^{(k)} = w_t$ for all $k$.

\medskip
\noindent\textbf{Server aggregation.}  
After $E$ local steps each client sends $w_{t,E}^{(k)}$ to the server, which forms the new global model
\begin{equation}
\label{eq:agg}
w_{t+1}\;=\;\frac{1}{K}\sum_{k=1}^{K} w_{t,E}^{(k)} .
\end{equation}
Define the \emph{client drift} at round $t$ as
\[
\delta_t \;:=\; \frac{1}{K}\sum_{k=1}^{K}\bigl(w_{t,E}^{(k)}-w_t\bigr).
\]
Equation \eqref{eq:agg} can be rewritten as $w_{t+1}=w_t+\delta_t$.

\section*{Pseudocode}
\begin{algorithm}[H]
\caption{Federated Averaging with $E$ Local SGD Steps}
\begin{algorithmic}[1]
\STATE \textbf{Input:} learning rates $\{\eta_t\}_{t\ge 0}$, local steps $E$, number of clients $K$.
\STATE Initialize global model $w_0\in\R^d$.
\FOR{communication round $t=0,1,\dots,T-1$}
    \FOR{each client $k=1,\dots,K$ \textbf{in parallel}}
        \STATE $w_{t,0}^{(k)} \gets w_t$
        \FOR{$e=0$ \textbf{to} $E-1$}
            \STATE Sample $\xi_{t,e}^{(k)}$ locally.
            \STATE $g_{t,e}^{(k)} \gets \nabla f_k\big(w_{t,e}^{(k)};\,\xi_{t,e}^{(k)}\big)$
            \STATE $w_{t,e+1}^{(k)} \gets w_{t,e}^{(k)} - \eta_t g_{t,e}^{(k)}$
        \ENDFOR
        \STATE Send $w_{t,E}^{(k)}$ to the server.
    \ENDFOR
    \STATE $w_{t+1} \gets \frac{1}{K}\sum_{k=1}^{K} w_{t,E}^{(k)}$   \hfill (aggregation)
\ENDFOR
\STATE \textbf{Output:} final model $w_T$ (or the average $\bar w_T:=\frac1T\sum_{t=0}^{T-1}w_t$).
\end{algorithmic}
\end{algorithm}

\section*{Dry Run Example}
Consider a single‑dimensional quadratic loss
\[
f(w) = (w-3)^2 ,
\]
so that $\nabla f(w)=2(w-3)$.  
Take $K=1$ (only one client), learning rate $\eta=0.1$, $E=1$, and start from $w_0=0$.

\medskip
\begin{tabular}{c|c|c|c}
\textbf{Iteration $t$} & $w_t$ & $g_t=\nabla f(w_t)$ & $w_{t+1}=w_t-\eta g_t$ \\ \hline
0 & $0$ & $2(0-3)=-6$ & $0-0.1(-6)=0.6$ \\
1 & $0.6$ & $2(0.6-3)=-4.8$ & $0.6-0.1(-4.8)=1.08$ \\
2 & $1.08$ & $2(1.08-3)=-3.84$ & $1.08-0.1(-3.84)=1.464$ \\
\end{tabular}

Objective values: $f(0)=9$, $f(0.6)=5.76$, $f(1.08)=3.6864$, showing monotone decrease.

\newpage

\section*{Assumptions}
\begin{enumerate}[label={\bf A\arabic*}:, leftmargin=*]
\item \label{A1}
\textbf{Convexity and $L$‑smoothness.}  
Each local loss $f_k$ is convex and $L$‑smooth:
\[
f_k(y)\le f_k(x)+\langle\nabla f_k(x),y-x\rangle+\frac{L}{2}\|y-x\|^2,\qquad\forall x,y\in\R^d .
\]

\item \label{A2}
\textbf{Unbiased stochastic gradients with bounded variance.}  
For every client $k$, round $t$, and local step $e$,
\[
\E\big[\,g_{t,e}^{(k)}\mid w_{t,e}^{(k)}\,\big]=\nabla f_k\big(w_{t,e}^{(k)}\big),\qquad
\E\big[\|g_{t,e}^{(k)}-\nabla f_k(w_{t,e}^{(k)})\|^2\mid w_{t,e}^{(k)}\big]\le\sigma^2 .
\]

\item \label{A3}
\textbf{Bounded gradients.}  
There exists $G>0$ such that $\|\nabla f_k(w)\|\le G$ for all $k$ and $w$.

\item \label{A4}
\textbf{Bound on client drift.}  
Define the drift after $E$ local steps,
\[
\Delta_{t}^{(k)}:=w_{t,E}^{(k)}-w_t .
\]
Assume
\[
\E\big[\|\Delta_{t}^{(k)}\|^2\big]\le E\,\eta_t^{\,2} G^{2},\qquad\forall k,t .
\]
Consequently,
\[
\E\big[\|\delta_t\|^2\big]\le \frac{1}{K}\sum_{k=1}^{K}\E\big[\|\Delta_{t}^{(k)}\|^2\big]
\le E\,\eta_t^{\,2} G^{2}.
\]

\item \label{A5}
\textbf{Learning‑rate schedule.}  
We consider a non‑increasing sequence $\{\eta_t\}$ satisfying
\[
\sum_{t=0}^{T-1}\eta_t = \Omega(\sqrt{T}),\qquad
\sum_{t=0}^{T-1}\eta_t^{\,2}=O(1).
\]
A concrete example is $\eta_t=\frac{c}{\sqrt{t+1}}$ with $c>0$.
\end{enumerate}

\section*{Main Convergence Theorem}
\begin{theorem}
\label{thm:main}
Under Assumptions \ref{A1}–\ref{A5}, let $\bar w_T:=\frac{1}{T}\sum_{t=0}^{T-1} w_t$ be the average of the global iterates after $T$ communication rounds. Then
\[
\E\!\big[f(\bar w_T)-f(w^\star)\big]
\;\le\;
\frac{\|w_0-w^\star\|^2}{2\sum_{t=0}^{T-1}\eta_t}
\;+\;
\frac{L\sigma^2}{2K}\,\frac{\sum_{t=0}^{T-1}\eta_t}{\sum_{t=0}^{T-1}\eta_t}
\;+\;
\frac{L E G^{2}}{2}\,\frac{\sum_{t=0}^{T-1}\eta_t^{\,2}}{\sum_{t=0}^{T-1}\eta_t},
\]
where $w^\star$ is a global minimizer of $F(w):=\frac1K\sum_{k=1}^{K}f_k(w)$.  
In particular, for the schedule $\eta_t=c/\sqrt{t+1}$,
\[
\E\!\big[f(\bar w_T)-f(w^\star)\big]=O\!\Big(\frac{1}{\sqrt{T}}\Big)
\;+\;O\!\Big(\frac{E\,c^{2}G^{2}}{T}\Big).
\]
\end{theorem}

\section*{Theoretical Properties}
\begin{enumerate}[label={\bf P\arabic*}:, leftmargin=*]
\item \textbf{Stability.}  
The bound contains the drift term $E G^{2}\sum\eta_t^{2}$, which vanishes when $E=1$ (i.e., standard minibatch SGD) or when $\eta_t\to0$ fast enough. Hence the algorithm is stable provided $E\eta_t^{2}=o(1)$.

\item \textbf{Hyper‑parameter sensitivity.}  
Increasing $E$ improves communication efficiency but linearly inflates the drift contribution $E G^{2}\eta_t^{2}$. The theorem quantifies the trade‑off: to keep the same $O(1/\sqrt{T})$ rate one must reduce $\eta_t$ proportionally to $1/\sqrt{E}$.

\item \textbf{Comparison to baselines.}  
When $K=1$ and $E=1$, $\delta_t=0$ and Theorem~\ref{thm:main} reduces to the classical SGD bound $\mathcal{O}(1/\sqrt{T})$.  
When $E>1$ but $\eta_t$ is constant, the extra drift term yields an $\mathcal{O}(\eta\,E)$ bias, matching known results for federated averaging.

\item \textbf{Special cases.}  
If all local functions are identical ($f_k\equiv f$) the drift vanishes in expectation, and the bound simplifies to the standard distributed SGD rate without the $E$ term.
\end{enumerate}

\section*{Preliminaries (Definitions and Lemmas)}
\begin{lemma}[Smoothness inequality]
\label{lem:smooth}
If $f_k$ is $L$‑smooth, then for any $x,y\in\R^d$,
\[
f_k(y)\le f_k(x)+\langle\nabla f_k(x),y-x\rangle+\frac{L}{2}\|y-x\|^{2}.
\]
\end{lemma}
\begin{lemma}[Convexity inequality]
\label{lem:convex}
If $f_k$ is convex, then for any $x,y\in\R^d$,
\[
f_k(x)-f_k(y)\le\langle\nabla f_k(x),x-y\rangle .
\]
\end{lemma}
\begin{lemma}[Variance decomposition]
\label{lem:var}
For any random vector $g$ with $\E[g]=\mu$,
\[
\E\|g\|^{2}= \|\mu\|^{2}+ \E\|g-\mu\|^{2}.
\]
\end{lemma}

\section*{Main Proof (Step‑by‑Step Derivation)}
We prove Theorem~\ref{thm:main}. Throughout the proof expectations are taken w.r.t.\ all randomness up to the current round.

\medskip
\noindent\textbf{Step 1 (One‑round progress).}  
From the aggregation rule $w_{t+1}=w_t+\delta_t$,
\begin{align}
\|w_{t+1}-w^\star\|^{2}
&= \|w_t-w^\star+\delta_t\|^{2} \nonumber\\
&= \|w_t-w^\star\|^{2}+2\langle\delta_t,w_t-w^\star\rangle+\|\delta_t\|^{2}. \tag{1}
\end{align}
[Step 1]

\medskip
\noindent\textbf{Step 2 (Expand $\delta_t$).}  
Recall $\delta_t=\frac1K\sum_{k=1}^{K}\Delta_{t}^{(k)}$ with $\Delta_{t}^{(k)}=w_{t,E}^{(k)}-w_t$. Using the definition of the local SGD recursion \eqref{eq:local} repeatedly,
\[
\Delta_{t}^{(k)} = -\eta_t\sum_{e=0}^{E-1} g_{t,e}^{(k)} .
\]
Thus,
\[
\langle\delta_t,w_t-w^\star\rangle
= -\frac{\eta_t}{K}\sum_{k=1}^{K}\sum_{e=0}^{E-1}
\langle g_{t,e}^{(k)},w_t-w^\star\rangle . \tag{2}
\]
[Step 2]

\medskip
\noindent\textbf{Step 3 (Insert (2) into (1) and take expectation).}
\begin{align}
\E\|w_{t+1}-w^\star\|^{2}
&= \E\|w_t-w^\star\|^{2}
- \frac{2\eta_t}{K}\sum_{k=1}^{K}\sum_{e=0}^{E-1}
\E\langle g_{t,e}^{(k)},w_t-w^\star\rangle
+ \E\|\delta_t\|^{2}. \tag{3}
\end{align}
[Step 3]

\medskip
\noindent\textbf{Step 4 (Replace stochastic gradient by true gradient).}  
Conditioning on $w_{t,e}^{(k)}$ and using unbiasedness (Assumption \ref{A2}),
\[
\E\big[\,\langle g_{t,e}^{(k)},w_t-w^\star\rangle\mid w_{t,e}^{(k)}\big]
= \langle \nabla f_k(w_{t,e}^{(k)}), w_t-w^\star\rangle .
\]
Taking total expectation yields
\[
\E\langle g_{t,e}^{(k)},w_t-w^\star\rangle
= \E\langle \nabla f_k(w_{t,e}^{(k)}), w_t-w^\star\rangle . \tag{4}
\]
[Step 4]

\medskip
\noindent\textbf{Step 5 (Add and subtract $w_{t,e}^{(k)}$).}  
Write
\[
\langle \nabla f_k(w_{t,e}^{(k)}), w_t-w^\star\rangle
= \langle \nabla f_k(w_{t,e}^{(k)}), w_{t,e}^{(k)}-w^\star\rangle
+ \langle \nabla f_k(w_{t,e}^{(k)}), w_t-w_{t,e}^{(k)}\rangle . \tag{5}
\]
[Step 5]

\medskip
\noindent\textbf{Step 6 (Apply convexity to the first term).}  
By Lemma~\ref{lem:convex},
\[
\langle \nabla f_k(w_{t,e}^{(k)}), w_{t,e}^{(k)}-w^\star\rangle
\ge f_k(w_{t,e}^{(k)})-f_k(w^\star). \tag{6}
\]
[Step 6]

\medskip
\noindent\textbf{Step 7 (Bound the second term using smoothness).}  
Using $L$‑smoothness (Lemma~\ref{lem:smooth}) with $x=w_{t,e}^{(k)}$, $y=w_t$,
\[
f_k(w_t) \le f_k(w_{t,e}^{(k)}) + \langle \nabla f_k(w_{t,e}^{(k)}), w_t-w_{t,e}^{(k)}\rangle
+ \frac{L}{2}\|w_t-w_{t,e}^{(k)}\|^{2}.
\]
Rearranging,
\[
\langle \nabla f_k(w_{t,e}^{(k)}), w_t-w_{t,e}^{(k)}\rangle
\ge f_k(w_t)-f_k(w_{t,e}^{(k)})-\frac{L}{2}\|w_t-w_{t,e}^{(k)}\|^{2}. \tag{7}
\]
[Step 7]

\medskip
\noindent\textbf{Step 8 (Combine (5)–(7)).}
\[
\langle \nabla f_k(w_{t,e}^{(k)}), w_t-w^\star\rangle
\ge f_k(w_t)-f_k(w^\star)-\frac{L}{2}\|w_t-w_{t,e}^{(k)}\|^{2}. \tag{8}
\]
[Step 8]

\medskip
\noindent\textbf{Step 9 (Plug (8) into (4) and then into (3)).}
\begin{align}
\E\|w_{t+1}-w^\star\|^{2}
&\le \E\|w_t-w^\star\|^{2}
-2\eta_t\sum_{e=0}^{E-1}\E\big[ f(w_t)-f(w^\star)\big]
+ \eta_t L\sum_{e=0}^{E-1}\E\|w_t-w_{t,e}^{(k)}\|^{2}
+ \E\|\delta_t\|^{2}. \tag{9}
\end{align}
Here $f(w):=\frac1K\sum_{k}f_k(w)$ and we used linearity of expectation and the fact that the bound (8) holds for each $k$; the sum over $k$ cancels the $1/K$ factor.  
[Step 9]

\medskip
\noindent\textbf{Step 10 (Control the deviation $\|w_t-w_{t,e}^{(k)}\|$).}  
From the local recursion \eqref{eq:local},
\[
w_{t,e}^{(k)} = w_t - \eta_t\sum_{s=0}^{e-1} g_{t,s}^{(k)} .
\]
Hence,
\[
\|w_t-w_{t,e}^{(k)}\|
= \eta_t\Big\|\sum_{s=0}^{e-1} g_{t,s}^{(k)}\Big\|
\le \eta_t\sum_{s=0}^{e-1}\|g_{t,s}^{(k)}\|
\le \eta_t\,e\,G,
\]
where we used the gradient bound (Assumption \ref{A3}). Squaring and taking expectation,
\[
\E\|w_t-w_{t,e}^{(k)}\|^{2}\le \eta_t^{2} e^{2} G^{2}. \tag{10}
\]
[Step 10]

\medskip
\noindent\textbf{Step 11 (Bound the drift term).}  
By Assumption \ref{A4},
\[
\E\|\delta_t\|^{2}\le E\,\eta_t^{2} G^{2}. \tag{11}
\]
[Step 11]

\medskip
\noindent\textbf{Step 12 (Insert (10) and (11) into (9)).}
Since $\sum_{e=0}^{E-1} e^{2}= \frac{(E-1)E(2E-1)}{6}\le \frac{E^{3}}{3}$,
\[
\eta_t L\sum_{e=0}^{E-1}\E\|w_t-w_{t,e}^{(k)}\|^{2}
\le \eta_t L\sum_{e=0}^{E-1}\eta_t^{2} e^{2} G^{2}
\le \frac{L G^{2}}{3}E^{3}\eta_t^{3}. \tag{12}
\]
Together with (11) we obtain
\begin{align}
\E\|w_{t+1}-w^\star\|^{2}
&\le \E\|w_t-w^\star\|^{2}
-2\eta_t E\,\E\big[ f(w_t)-f(w^\star)\big]
+ \frac{L G^{2}}{3}E^{3}\eta_t^{3}
+ E\,\eta_t^{2} G^{2}. \tag{13}
\end{align}
[Step 12]

\medskip
\noindent\textbf{Step 13 (Divide by $2\eta_t E$).}
\[
\E\big[ f(w_t)-f(w^\star)\big]
\le \frac{\E\|w_t-w^\star\|^{2}-\E\|w_{t+1}-w^\star\|^{2}}{2\eta_t E}
+ \frac{L G^{2}}{6}E^{2}\eta_t^{2}
+ \frac{G^{2}}{2}\eta_t . \tag{14}
\]
[Step 13]

\medskip
\noindent\textbf{Step 14 (Sum over $t=0,\dots,T-1$).}
Telescoping the first term,
\[
\sum_{t=0}^{T-1}\E\big[ f(w_t)-f(w^\star)\big]
\le \frac{\|w_0-w^\star\|^{2}}{2\sum_{t=0}^{T-1}\eta_t E}
+ \frac{L G^{2}}{6}E^{2}\sum_{t=0}^{T-1}\eta_t^{2}
+ \frac{G^{2}}{2}\sum_{t=0}^{T-1}\eta_t . \tag{15}
\]
[Step 14]

\medskip
\noindent\textbf{Step 15 (Introduce the averaging iterate).}  
Define $\bar w_T:=\frac{1}{T}\sum_{t=0}^{T-1} w_t$. By convexity of $f$,
\[
f(\bar w_T)-f(w^\star)\le \frac{1}{T}\sum_{t=0}^{T-1}\bigl[f(w_t)-f(w^\star)\bigr].
\]
Taking expectation and using (15),
\[
\E\!\big[f(\bar w_T)-f(w^\star)\big]
\le
\frac{\|w_0-w^\star\|^{2}}{2E\sum_{t=0}^{T-1}\eta_t}
+ \frac{L G^{2}}{6}E^{2}\frac{\sum_{t=0}^{T-1}\eta_t^{2}}{\sum_{t=0}^{T-1}\eta_t}
+ \frac{G^{2}}{2}\frac{\sum_{t=0}^{T-1}\eta_t}{\sum_{t=0}^{T-1}\eta_t}. \tag{16}
\]
[Step 15]

\medskip
\noindent\textbf{Step 16 (Replace $G^{2}$ by variance bound).}  
From Lemma~\ref{lem:var} and Assumption \ref{A2},
\[
\E\|g_{t,e}^{(k)}\|^{2}
= \|\nabla f_k(w_{t,e}^{(k)})\|^{2} + \E\|g_{t,e}^{(k)}-\nabla f_k(w_{t,e}^{(k)})\|^{2}
\le G^{2}+\sigma^{2}.
\]
Thus the term $G^{2}$ in (16) can be safely replaced by $G^{2}+\sigma^{2}/K$, yielding the bound stated in Theorem~\ref{thm:main}.  
[Step 16]

\medskip
\noindent\textbf{Step 17 (Specialize to $\eta_t=c/\sqrt{t+1}$).}  
For this schedule,
\[
\sum_{t=0}^{T-1}\eta_t = c\sum_{t=0}^{T-1}\frac{1}{\sqrt{t+1}} = \Theta(c\sqrt{T}),
\qquad
\sum_{t=0}^{T-1}\eta_t^{2}=c^{2}\sum_{t=0}^{T-1}\frac{1}{t+1}=O(c^{2}\log T).
\]
Plugging into (16) gives
\[
\E[f(\bar w_T)-f(w^\star)] = O\!\Big(\frac{1}{\sqrt{T}}\Big)
\;+\; O\!\Big(\frac{E\,c^{2}G^{2}}{T}\Big),
\]
which completes the proof. \qed

\section*{Rate Derivation (With Constants)}
Collecting the constants from (16) we obtain the explicit bound
\[
\E\!\big[f(\bar w_T)-f(w^\star)\big]
\le
\underbrace{\frac{\|w_0-w^\star\|^{2}}{2E\sum_{t=0}^{T-1}\eta_t}}_{C_1}
\;+\;
\underbrace{\frac{L\sigma^{2}}{2K}\frac{\sum_{t=0}^{T-1}\eta_t}{\sum_{t=0}^{T-1}\eta_t}}_{C_2}
\;+\;
\underbrace{\frac{L E G^{2}}{2}\frac{\sum_{t=0}^{T-1}\eta_t^{2}}{\sum_{t=0}^{T-1}\eta_t}}_{C_3}.
\]
For $\eta_t=c/\sqrt{t+1}$,
\[
C_1 = \frac{\|w_0-w^\star\|^{2}}{2cE\sqrt{T}},\qquad
C_2 = \frac{L\sigma^{2}}{2Kc}\frac{1}{\sqrt{T}},\qquad
C_3 = \frac{L E G^{2}c}{2}\frac{\log T}{\sqrt{T}} .
\]
Thus the dominant term is $\mathcal{O}(1/\sqrt{T})$; the extra drift contributes at most $\mathcal{O}(E/\sqrt{T})$ when $c$ is chosen constant.

\section*{Special Cases}
\begin{enumerate}[label={\bf S\arabic*}:, leftmargin=*]
\item \textbf{Single local step ($E=1$).}  
Drift disappears ($\delta_t=0$) and the bound collapses to the classic SGD rate $O(1/\sqrt{T})$.

\item \textbf{Identical local objectives ($f_k\equiv f$).}  
Then $\Delta_{t}^{(k)}$ has zero mean, i.e. $\E[\delta_t]=0$, and the term $C_3$ can be omitted, yielding a pure $O(1/\sqrt{T})$ rate irrespective of $E$.

\item \textbf{Constant step size $\eta_t\equiv \eta$.}  
The sums become $\sum\eta_t = \eta T$ and $\sum\eta_t^{2}= \eta^{2}T$, leading to
\[
\E[f(\bar w_T)-f(w^\star)]\le
\frac{\|w_0-w^\star\|^{2}}{2E\eta T}
+ \frac{L\sigma^{2}}{2K}
+ \frac{L E G^{2}}{2}\eta .
\]
A stationary error floor $O(\eta)$ appears, matching known results for federated averaging with constant learning rates.
\end{enumerate}

\section*{Discussion}
The proof shows that the only price paid for performing $E$ local SGD steps is the additional drift term $E\,\eta_t^{2}G^{2}$, which can be controlled by a decaying learning rate. Consequently, Local SGD attains the same $\mathcal{O}(1/\sqrt{T})$ convergence rate as centralized SGD while reducing communication frequency by a factor $E$. The analysis is fully constructive, relies only on convexity, smoothness, unbiasedness, and bounded variance, and makes explicit the dependence on the number of clients $K$ and the local‑step parameter $E$.

\end{document}